%!TEX program = xelatex
\documentclass[a4paper,12pt]{article}
\usepackage{fontspec}\defaultfontfeatures{Ligatures=TeX}
\usepackage[a4paper,vmargin={3cm,3.5cm},hmargin={3cm,3cm}]{geometry}
%-----------------------------------------------------------------------------%
\usepackage{settings}
%-----------------------------------------------------------------------------%
%%% Title %%%
\title{G6411 Recitation Note: \glsentryshort{ols} Algebraic Properties}
\author{Jesse Chieh Chen\\\href{mailto:cc5229@columbia.edu}{cc5229@columbia.edu}}
\date{\today}
%-----------------------------------------------------------------------------%

\begin{document}

% We need to cover:
%
% # Define what OLS is
%
% # Two approaches to OLS
%
% two approaches leading to the same estimator, and then introduce the projection interpretation of OLS, both sample and popualtion
%
% 1. Loss function approach
%     - Show clearly the sample/population rank condition for identification
% 2. Analogue approach: Method of Moments
%     - Conditional mean restriction: interpret this in terms of both sample and population
%
% # Projection interpretation
%
% Introduce the projection and annihilation matrices, and show how they are related to the OLS estimator.
% Then introduce the geometric interpretation of OLS.
%
% # FWL Theorem
%
% First define and proof and explain, then introduce the graphical interpretation of FWL theorem.
%
% # Goodness of fit
%
% - R squred: Define, explain, and interpret geometrically (Pythagorean theorem)
% - Adjusted R squared: Define, explain, and interpret geometrically (Pythagorean theorem)
% - relation to cosine similarity: Define, explain, and interpret geometrically (Pythagorean theorem)

\maketitle

\noindent
% In practice, we do not know the population distribution of $(Y_i,X_i)$, so we cannot compute the population least squares coefficient.
We observe a sample $\{(Y_i,X_i)\}_{i=1}^n$,
where $Y_i$ is a scalar and $X_i=(X_{i1},\ldots,X_{ik})\T$ is a vector of regressors.
We stack the outcomes as $Y=(Y_1,\ldots,Y_n)\T\in\reals^n$
and the regressors as $X=[X_1\;\cdots\;X_n]\T\in\reals^{n\times k}$,
so the $i$-th row of $X$ is $X_i\T$.
The matrix $X$ is sometimes known as the \emph{design matrix}.
The algebraic results below hold for any observed sample:
they do not require \gls{iid} sampling or a correctly specified conditional mean.

\section{Two Approaches to \glsentryshort{ols}}

\begin{definition}[\glsentrylong{ols}]\label{def-ols}
	The \gls{ols} estimator is a solution to
	\begin{align}\label{eq-ols-objective}
		\widehat\beta\in\arg\min_{b\in\reals^k}\frac{1}{n}\sum_{i=1}^n(Y_i-X_i\T b)^2.
	\end{align}
	The objective \eqref{eq-ols-objective} can also be written as $n\inv(Y-X b)\T(Y-X b)$.
\end{definition}

\subsection{The Loss Function Approach}

In this approach, we directly solve for the minimizer of the sample least squares objective \eqref{eq-ols-objective}.

% To find the \gls{ols} coefficient to minimize \eqref{eq-ols-objective},
% we can directly solve for the first order condition:
% \begin{align}\label{eq-sample-first-order-condition}
% 	\frac{1}{n}\sum_{i=1}^n X_i(Y_i-X_i\T\widehat\beta)=0.
% \end{align}

\begin{theorem}[\gls{ols} Estimator]\label{thm-sample-least-squares}
	Suppose that $X$ has full column rank $k$.
	Then the \gls{ols} estimator is unique and equals
	\begin{align}\label{eq-ols-formula}
		\widehat\beta
		&=(X\T X)\inv X\T Y\\
		&=\left(\frac{1}{n}\sum_{i=1}^nX_iX_i\T\right)\inv
		\left(\frac{1}{n}\sum_{i=1}^nX_iY_i\right).
	\end{align}
	Equivalently, it is the unique solution to the equations
	\begin{align}\label{eq-sample-normal-equations}
		X\T(Y-X\widehat\beta)=0.
	\end{align}
\end{theorem}
\begin{proof}
	The gradient of the sample objective is
	\begin{align}
		-\frac2n \sum_{i=1}^n X_iY_i + \frac2n \sum_{i=1}^n X_iX_i\T b
		=-\frac2n X\T Y + \frac2n X\T X b.
	\end{align}
	Setting it equal to zero gives
	\begin{align}
		\frac1n \sum_{i=1}^{n} X_i(Y_i-X_i\T\widehat\beta)
		=\frac1n X\T(Y-X\widehat\beta)=0.
	\end{align}
	This yields the desired solution
	\begin{align}
		\begin{aligned}
			\widehat\beta
			=\left(\frac1n\sum_{i=1}^{n}X_iX_i\T\right)\inv
			\left(\frac1n\sum_{i=1}^{n}X_iY_i\right)
			=(X\T X)\inv X\T Y.
		\end{aligned}
	\end{align}
	Full column rank makes $X\T X$ positive definite,
	so the objective is strictly convex and its unique minimizer is given by equation~\eqref{eq-ols-formula}.
\end{proof}

\begin{remark}[The Sample Counterpart of Population Regression]
	For population statements, assume that $Y_i$ and the components of $X_i$ have finite second moments.
	The population regression coefficient minimizes expected squared loss:
	\begin{align}
		\beta\in\arg\min_{b\in\reals^k}\E[(Y_i-X_i\T b)^2],
	\end{align}
	The sample objective \eqref{eq-ols-objective} replaces this expectation by an average over the observed data.
\end{remark}

\begin{remark}[Population Identification and Sample Uniqueness]
	In the population, $Q\coloneqq\E(X_iX_i\T)$ being nonsingular identifies a unique coefficient $\beta$.
	In the sample, full column rank of $X$ makes $X\T X$ nonsingular and gives a unique estimate $\widehat\beta$.
	This requires $n\ge k$ and no exact linear dependence among the columns of $X$.
	The two conditions are different: even when $Q$ is nonsingular,
	a particular sample may lack enough variation to give full column rank.
	If $Xv=0$ for some nonzero $v$, then $b$ and $b+v$ give the same fitted vector.
\end{remark}

\subsection{The Analogue Approach: Method of Moments}

The \emph{method of moments} replaces population moment equations by sample moment equations.
The population least squares coefficient satisfies the orthogonality conditions
\begin{align}
	\E(X_i u_i)=\E[X_i(Y_i-X_i\T\beta)]=0
	\quad\text{where}\quad
	u_i\coloneqq Y_i-X_i\T\beta.
\end{align}
Instead of minimizing a loss function, we can start from these equations
and replace the expectation by a sample average:
\begin{align}\label{eq-sample-orthogonality}
	\frac{1}{n}\sum_{i=1}^n X_i(Y_i-X_i\T b)=0.
\end{align}
These are exactly the sample normal equations.
Under full column rank, solving them gives the same \gls{ols} estimator as the loss function approach.

\begin{remark}[Conditional Mean versus Orthogonality]
	A model may impose the stronger restriction $\E(u_i\given X_i)=0$:
	the error averages to zero at each regressor value, so $\E(Y_i\given X_i)=X_i\T\beta$.
	This implies $\E(X_i u_i)=0$, but the converse need not hold.
	The population least squares coefficient satisfies orthogonality even when the conditional mean is nonlinear.
	In the sample, \gls{ols} sets $n\inv\sum_i X_i\widehat u_i=0$
	where $\widehat u_i=Y_i-X_i\T\widehat\beta$ by construction.
	Similar to the population regression,
	it does not generally set the average residual to zero separately at each regressor value.
\end{remark}

\section{Projection Interpretation}

\begin{figure}[t]
	\centering
	\includegraphics[scale=1]{figures/ols.pdf}
	\caption{The Geometry of \gls{ols}.}
	\label{fig-ols}
\end{figure}

For the matrix formulae that follow, assume that $X$ has full column rank.
Write $I_n$ for the $n\times n$ identity matrix and $\|a\|^2=a\T a$ for squared Euclidean length.

\begin{definition}[Fitted Values and Residuals]\label{def-fitted-values-residuals}
	The fitted value for observation $i$ is $\widehat Y_i\coloneqq X_i\T\widehat\beta$,
	and its residual is $\widehat u_i\coloneqq Y_i-\widehat Y_i$.
	In vector form,
	\begin{align}
		\widehat Y=X\widehat\beta,
		\qquad \widehat u\coloneqq(\widehat u_1,\ldots,\widehat u_n)\T=Y-\widehat Y.
	\end{align}
\end{definition}

\begin{proposition}[Sample Orthogonality]\label{prop-sample-orthogonality}
	The \gls{ols} residuals satisfy
	\begin{align}
		X\T\widehat u=0
		\quad\text{and}\quad
		\widehat Y\T\widehat u=0.
	\end{align}
	If an intercept is included, then
	\begin{align}
		\sum_{i=1}^n\widehat u_i=0,
	\end{align}
	and the sample average of the fitted values equals the sample average of the outcomes.
\end{proposition}
\begin{proof}
	The first equality follows from $\widehat\beta$ being the solution to \eqref{eq-sample-orthogonality}:
	\begin{align}
		X\T\widehat u
		=X\T(Y-X\widehat\beta)
		=X\T Y-X\T X(X\T X)\inv X\T Y
		=X\T Y-X\T Y
		=0.
	\end{align}
	The second equality follows by substituting $\widehat Y=X\widehat\beta$ and using the first equality:
	\begin{align}
		\widehat Y\T\widehat u
		=(X\widehat\beta)\T\widehat u
		=\widehat\beta\T(X\T\widehat u)
		=0.
	\end{align}
	For the third equality, suppose that the first column of $X$ is $\one$, the vector of ones.
	Then the first row of $X\T\widehat u=0$ gives $\one\T\widehat u=0$.
	Since $\widehat u_i=Y_i-\widehat Y_i$, this also gives $n\inv\sum_i\widehat Y_i=n\inv\sum_iY_i$.
\end{proof}

\begin{definition}[Projection and Annihilation Matrices]\label{def-projection-matrices}
	The \emph{projection matrix} and \emph{annihilation matrix} are
	\begin{align}
		P\coloneqq X(X\T X)\inv X\T,
		\qquad M\coloneqq I_n-P.
	\end{align}
	They give the fitted values and residuals directly: $\widehat Y=PY$ and $\widehat u=MY$.
	The matrix $M$ is also called the \emph{residual-maker matrix}.
\end{definition}

% \begin{remark}[What is a projection?]
% 	A projection is a linear mapping such that when applied twice, it gives the same result as when applied once.
% 	For a matrix $A$, this means $A^2=A$, a property called \emph{idempotence}.
% 	When $A$ is also symmetric, it is an \emph{orthogonal} projection:
% 	the part removed is perpendicular to the space onto which we project.
% \end{remark}

\begin{proposition}[Properties of $P$ and $M$]\label{prop-projection-matrices}
	Both matrices are symmetric and idempotent:
	\begin{align}
		P\T=P,\quad P^2=P,
		\qquad M\T=M,\quad M^2=M.
	\end{align}
	Furthermore,
	\begin{align}
		PX=X,\qquad MX=0,\qquad PM=MP=0.
	\end{align}
\end{proposition}
\begin{proof}
	Symmetry follows by transposing $P$ and $M$.
	For idempotence,
	\begin{align*}
		P^2&=X(X\T X)\inv(X\T X)(X\T X)\inv X\T=P,\\
		M^2&=(I_n-P)^2=I_n-2P+P^2=M.
	\end{align*}
	Also, $PX=X(X\T X)\inv X\T X=X$, so $MX=X-PX=0$.
	Finally, $PM=P-P^2=0$, and similarly $MP=0$.
\end{proof}

\begin{remark}[Projection Interpretation of \gls{ols}]
	Consider \autoref{fig-ols}.
	Here the labels $X_1$ and $X_2$ represent two columns of the design matrix, rather than individual observations.
	The fitted vector $\widehat Y$ or $PY$ lies in the column space of $X$,
	and the residual vector is orthogonal to that space.
	Thus, \gls{ols} projects $Y$ onto the space of vectors that can be written as $X b$,
	i.e., the span of the columns of $X$.
	Applying $P$ again does not change the fitted vector; applying $M$ to any vector $Xb$ removes it completely.
	The decomposition $Y=PY+MY$ forms a right triangle, so the Pythagorean theorem gives
	\begin{align}\label{eq-uncentered-pythagoras}
		\|Y\|^2=\|PY\|^2+\|MY\|^2.
	\end{align}
\end{remark}

\begin{remark}[The Population Counterpart]
	Let $u_i=Y_i-X_i\T\beta$ denote the population regression error.
	In the sample, orthogonality means that the usual dot product is zero.
	In the population, we instead require the expected product to be zero.
	The condition $\E(X_i u_i)=0$ implies that $u_i$ is orthogonal to every linear combination of the regressors:
	\begin{align}
		\E[(X_i\T b)u_i]=b\T\E(X_i u_i)=0
		\quad\forall b\in\reals^k.
	\end{align}
	Thus, $X_i\T\beta$ is the population projection of $Y_i$ onto linear functions of $X_i$.
	Taking $b=\beta$ eliminates the cross term and gives the population version of the Pythagorean identity:
	\begin{align}
		\E[Y_i^2]
		= \E[(X_i\T\beta + u_i)^2]
		= \E[(X_i\T\beta)^2] + \E[u_i^2].
	\end{align}
\end{remark}

\section{The \glsentrylong{fwl} Theorem}

Partition the design matrix into two groups of columns, $X=[X_1\;X_2]$,
with $k_1$ and $k_2$ columns, respectively, where $k_1+k_2=k$.
In this section, $X_1$ and $X_2$ denote \emph{blocks of the design matrix}, not individual observations.
Write $\widehat\beta=(\widehat\beta_1\T,\widehat\beta_2\T)\T$ for the corresponding coefficient blocks.
If the regression includes an intercept, place it in $X_1$.
Define
\begin{align}
	P_1=X_1(X_1\T X_1)\inv X_1\T,
	\qquad M_1=I_n-P_1.
\end{align}

\begin{theorem}[\gls{fwl}]\label{thm-fwl}
	Suppose that $[X_1\;X_2]$ has full column rank.
	The coefficient on $X_2$ in the regression of $Y$ on $[X_1\;X_2]$ is
	\begin{align}\label{eq-fwl}
		\widehat\beta_2=(X_2\T M_1X_2)\inv X_2\T M_1Y.
	\end{align}
	Equivalently, regress $M_1Y$ on $M_1X_2$ without adding an intercept.
	This regression gives the same coefficient $\widehat\beta_2$ and the same residual vector as the full regression.
\end{theorem}
\begin{proof}
	The first block of the full normal equations gives
	\begin{align*}
		\widehat\beta_1=(X_1\T X_1)\inv X_1\T(Y-X_2\widehat\beta_2).
	\end{align*}
	Substituting this into the full residual gives
	\begin{align*}
		\widehat u=Y-X_1\widehat\beta_1-X_2\widehat\beta_2
		=M_1Y-M_1X_2\widehat\beta_2.
	\end{align*}
	The second block of the normal equations is therefore
	\begin{align*}
		X_2\T M_1Y-X_2\T M_1X_2\widehat\beta_2=0.
	\end{align*}
	The matrix $M_1X_2$ has full column rank:
	otherwise, a nonzero combination of the columns of $X_2$ would lie in the span of $X_1$,
	contradicting full column rank of $[X_1\;X_2]$.
	Since $M_1\T M_1=M_1$, the cross-product matrix is
	$(M_1X_2)\T(M_1X_2)=X_2\T M_1X_2$, which is nonsingular.
	Solving gives \eqref{eq-fwl}, and the displayed residual identity proves the remaining claim.
\end{proof}

\begin{remark}[What Does ``Controlling For'' Mean?]
	The coefficient on $X_2$ is often described as ``its coefficient after controlling for $X_1$.''
	What does ``controlling for'' mean in this context?
	The theorem gives a three-step recipe:
	\begin{enumerate}
		\item Regress $Y$ on $X_1$ and keep the residual $M_1Y$.
		\item Regress every column of $X_2$ on $X_1$ and keep the residuals $M_1X_2$.
		\item Regress $M_1Y$ on $M_1X_2$.
	\end{enumerate}
	Thus, the coefficient on $X_2$ uses the variation left after removing the linear contribution of $X_1$ from both the outcome and the regressors.
\end{remark}

\begin{example}[Demeaning as a Special Case]\label{eg-fwl-demeaning}
	If $X_1=\one$, the vector of ones, then $P_1=n\inv\one\one\T$.
	Multiplying by $M_1$ subtracts each variable's sample mean.
	Thus, regressing an outcome on regressors and an intercept gives the same slope coefficients
	as regressing the demeaned outcome on the demeaned regressors without an intercept.
\end{example}

\begin{remark}[The Geometry of \glsentryshort{fwl}]
	In \autoref{fig-fwl}, each block contains a single regressor vector.
	The full regression projects $Y$ onto the plane spanned by $X_1$ and $X_2$.
	Removing the $X_1$ component leaves the red vector $M_1Y$ and the blue direction $M_1X_2$.
	Projecting the red vector onto that direction gives $M_1X_2\widehat\beta_2$.
	The remaining perpendicular component is the same residual $Y-\widehat Y$ as in the full regression.
	Note that vectors are translated in the drawing to show these right triangles together.
\end{remark}

\begin{figure}[t]
	\centering
	\includegraphics[scale=1]{figures/fwl.pdf}
	\caption{The Geometry of \gls{fwl} Theorem.}
	\label{fig-fwl}
\end{figure}

\section{Goodness of Fit}

\subsection{The Share Explained by the Projection}

Recall that the fitted values and residuals give the orthogonal decomposition $Y=PY+MY$.
The Pythagorean theorem therefore gives
\begin{align}
	\underbrace{\|Y\|^2}_{\text{total squared length}}
	=\underbrace{\|PY\|^2}_{\text{explained by the projection}}
	+\underbrace{\|MY\|^2}_{\text{unexplained}}.
\end{align}
A measure of fit is the share of the total squared length accounted for by $PY$.

\begin{definition}[Uncentered Coefficient of Determination]
	For $Y\ne0$, the \emph{uncentered} $R^2$ is
	\begin{align}\label{eq-uncentered-r-squared}
		R_u^2\coloneqq\frac{\|PY\|^2}{\|Y\|^2}
		=1-\frac{\|MY\|^2}{\|Y\|^2}
		=1-\frac{\sum_{i=1}^n\widehat u_i^2}{\sum_{i=1}^nY_i^2}.
	\end{align}
\end{definition}

\begin{remark}
	The orthogonal decomposition ensures that $0\le R_u^2\le1$, with or without an intercept.
	However, this measure counts the entire fitted vector $PY$ as explained,
	including any part that simply predicts the mean.
\end{remark}

\begin{example}[Explaining Mostly the Mean]\label{eg-r-squared-mean}
	Let $Y=(9,11)\T$ and regress it on an intercept alone.
	Then $PY=(10,10)\T$ and $MY=(-1,1)\T$, so
	\begin{align}
		R_u^2=\frac{10^2+10^2}{9^2+11^2}=\frac{100}{101}\approx0.99.
	\end{align}
	The fitted vector accounts for almost all of $\|Y\|^2$,
	yet it explains none of the differences between observations: it predicts the same mean for everyone.
\end{example}

\subsection{Partialling Out the Mean: Centered \texorpdfstring{$R^2$}{R-squared}}

To focus on variation \emph{around the mean}, we first remove the mean using \gls{fwl}.
Suppose that the regression includes an intercept, and write $X=[\one\;X_2]$.
We have
\begin{align}
	P_1=\frac1n\one\one\T,
	\qquad M_1=I_n-\frac1n\one\one\T.
\end{align}
As explained in \autoref{eg-fwl-demeaning},
partialling out the intercept centers the variables.
By \gls{fwl}, regressing $M_1 Y$ on $M_1 X_2$ gives the same slope coefficients and residuals as the original regression.
Here $P$ and $M$ continue to denote the projection and residual-maker matrices for the full design matrix $X$.
Let $\bar Y=n\inv\sum_{i=1}^nY_i$.
Since the regression includes an intercept, the fitted values average to $\bar Y$ and the residuals average to zero.
Thus, $M_1PY=PY-\bar Y\one$ and $M_1MY=MY$.
Applying $M_1$ to $Y=PY+MY$ gives
\begin{align}
	M_1Y=M_1PY+MY.
\end{align}
% The fitted vector in the residualized regression is $M_1X_2\widehat\beta_2=M_1PY$,
% and its residual is $MY$.
These are orthogonal, just as fitted values and residuals are in the original regression.

\begin{proposition}[Decomposition of Variation]\label{prop-centered-pythagoras}
	The centered total, explained, and residual sums of squares satisfy
	\begin{align}\label{eq-centered-pythagoras}
		\underbrace{\|M_1 Y\|^2}_{\mathrm{TSS}}
		=\underbrace{\|M_1PY\|^2}_{\mathrm{ESS}}
		+\underbrace{\|MY\|^2}_{\mathrm{RSS}}.
	\end{align}
\end{proposition}
\begin{proof}
	Note the following:
	\begin{align}
		\|M_1 Y\|^2
		= \|M_1PY\|^2 + \|MY\|^2 + 2 Y\T P M_1 MY.
	\end{align}
	Since $P M_1 M = P M = 0$, we have the desired decomposition.
\end{proof}

\begin{definition}[Centered Coefficient of Determination]
	For nonconstant $Y$ and a regression including an intercept,
	the \emph{centered} $R^2$, usually simply called $R^2$, is
	\begin{align}\label{eq-centered-r-squared}
		R^2\coloneqq\frac{\|M_1PY\|^2}{\|M_1 Y\|^2}
		=\frac{\mathrm{ESS}}{\mathrm{TSS}}
		=1-\frac{\mathrm{RSS}}{\mathrm{TSS}}.
	\end{align}
\end{definition}

\begin{remark}[Reading Centered $R^2$]
	$R^2$ is the fraction of the sample variation around $\bar Y$ accounted for by the centered fitted values.
	The idea is the same as for uncentered $R_u^2$, but we measure the projection's share \emph{after removing the mean}.
	In \autoref{eg-r-squared-mean}, $M_1PY=0$, so centered $R^2=0$.
\end{remark}

\begin{remark}[Why the Intercept Matters]
	Without an intercept, the residuals need not sum to zero, and centering the outcome and regressors generally changes the fitted regression.
	Consequently, the centered decomposition need not hold for the original fit.
	If we still report
	\begin{align}
		1-\frac{\sum_{i=1}^n\widehat u_i^2}{\sum_{i=1}^n(Y_i-\bar Y)^2},
	\end{align}
	it can be negative and need not equal a share of explained variation.
	A negative value means that the regression fits worse than simply predicting $\bar Y$.
	In contrast, the uncentered decomposition $Y=PY+MY$ remains orthogonal without an intercept,
	so $R_u^2$ remains between zero and one whenever it is defined.
\end{remark}

\subsection{Adjusted \texorpdfstring{$R^2$}{R-squared}}

\begin{definition}[Adjusted Coefficient of Determination]
	Suppose that the regression includes an intercept and $n>k$,
	where $k$ counts all coefficients, including the intercept.
	The \emph{adjusted} version of centered $R^2$ is
	\begin{align}
		\bar R^2\coloneqq1-\frac{\mathrm{RSS}/(n-k)}{\mathrm{TSS}/(n-1)}
		=1-\frac{n-1}{n-k}(1-R^2).
	\end{align}
\end{definition}

\begin{remark}
	Adding regressors on the same sample cannot decrease $R^2$, but can decrease $\bar R^2$.
	When a new regressor does not improve the fit enough to offset the loss of degrees of freedom, adjusted $R^2$ decreses.
\end{remark}

\begin{remark}[Why Divide by $n-k$ and $n-1$?]
	One statistical explanation uses variance estimation.
	For \gls{iid} observations with a linear conditional mean and homoskedastic error with variance $\sigma^2$,
	\begin{align*}
		\E(\mathrm{RSS}\given X)&=(n-k)\sigma^2,\\
		\E(\mathrm{TSS})&=(n-1)\var(Y_i).
	\end{align*}
	Thus, $\mathrm{RSS}/(n-k)$ and $\mathrm{TSS}/(n-1)$ are unbiased estimates of the error variance and outcome variance, respectively.
	Adjusted $R^2$ compares these corrected variance estimates.
	This does not make their ratio, or adjusted $R^2$ itself, unbiased.
\end{remark}

\vfill
\section*{Acronyms}
\printglossaries

\end{document}
